\documentclass[11pt,a4paper]{article}
\usepackage[margin=2.6cm]{geometry}
\usepackage{amsmath,amssymb}
\usepackage[authoryear,round]{natbib}
\usepackage{xcolor}
\usepackage[colorlinks=true,linkcolor=blue!50!black,citecolor=blue!45!black]{hyperref}

\newcommand{\PB}{P^{\mathrm{B}}}
\newcommand{\PJ}{P^{\mathrm{J}}}
\newcommand{\assoc}{\operatorname{assoc}}
\newcommand{\bigO}{O}

\title{The Paper's Dialectic\\
\large A positioning document for \emph{Order-Sensitivity of Belief and Decision
Statistics under Jeffrey Conditioning}}
\author{}
\date{}

\begin{document}
\maketitle
\vspace{-2.4em}

\begin{center}\small\itshape
Standalone working document, not manuscript content. Every quotation below was
read in the primary source, and each carries its page. The one result cited
without direct reading is Jeffrey (1988), whose status is stated in the
references.
\end{center}

\section*{The standoff}

The case against sequential Jeffrey updating is a psychological claim made by
people who are not psychologists. Its sharpest statement is the title of
\citet{Doring1999}: \emph{Why Bayesian Psychology Is Incomplete}. His abstract
carries the whole indictment:

\begin{quote}
``Jeffrey conditionalization is sensitive to the order in which the evidence
arrives. This order effect can be so pronounced as to call for a belief
adjustment that cannot be understood as an assimilation of incoming evidence by
Jeffrey's rule.'' (S379)
\end{quote}

\citet{Hawthorne2004}, who named the model, agrees with the verdict while
disclaiming the competence to deliver it, in the same breath:

\begin{quote}
``Which extension of Basic Jeffrey Updating is the more plausible model of human
agents? I'm a logician, not a psychologist. But Amnestic Updating seems
psychologically less plausible than the more Bayesian approaches, not because it
is un-Bayesian, but because it seems unlikely that we dismiss previous
experiences so completely.'' (p.~115)
\end{quote}

That is the standoff. The attack says the model is unrealistic. The attack is
about psychology. Nobody in the exchange measures anything. The paper's
contribution is to say what measurement could and could not settle here.

\section*{The dialectic in four steps}

\paragraph{1. For soft evidence, the order effect is the price of coherence.}
An impression leaves no certifiable proposition, so strict conditionalisation is
unavailable. If the impression delivers a level, a revised credence on the
attribute's partition, then Jeffrey updating is the unique coherent revision and
the minimal change of the prior consistent with it \citep{DiaconisZabell1982}.
Applied to a second cue it holds the last-read marginal exactly, and
sequence-dependence follows. The order effect is not a defect bolted onto the
rule. It is what coherence costs when inputs arrive as levels.

\paragraph{2. The escape re-describes the input, and even Jeffrey took it.}
\citet{Field1978} reparametrised the update so that the input is a portable
factor rather than a delivered credence, and sequential revisions then commute.
\citet{Jeffrey1988} proved a commutativity result of the same kind for
probability factors. \citet{Wagner2002} completed the line: the
non-commutativity ``has caused much unjustified concern,'' because ``when
identical learning is properly represented, namely, by identical Bayes factors
rather than identical posterior probabilities, then sequential
probability-kinematical revisions behave just as they should'' (abstract). The escape is
blocked twice over. Phenomenologically, a Bayes factor requires the probability
of the same credential for a candidate who is not competent, and an impression
supplies no such counterfactual; the panelist feels a level, not a ratio.
Internally, \citet{Garber1980} showed Field's portable parameter compounds
absurdly under repetition, and even the attacker rejects the escape:
``Field has neither shown that Jeffrey conditionalization is inapplicable nor
that his own scheme is a viable alternative'' \citep[S388]{Doring1999}.

\paragraph{3. The attack has two prongs, and the model here inherits both.}
Hawthorne's first prong is that the overwriting is implausible in itself: it
``seems implausible that the most recent experience or non-propositional state
should completely dictate belief strengths for basis sentences, with no regard
for the import of previous experiences or states'' (p.~99). His illustration is a
single basis cued twice---``when I next see my car in dim light, its appearance
at that moment completely overwrites my beliefs regarding its true color''
(p.~97)---and in that form it does not reach a model whose cues address one
attribute each. The complaint itself does reach it. A credential that shifts the
panel's view of trustworthiness through the believed correlation has that
implication erased when the letter arrives and sets trustworthiness to its own
impression: the import of the earlier experience is indeed disregarded. The
second prong is about consequences: ``the order in which states or experiences
are acquired will almost always have a very significant influence on an agent's
belief strengths. This order effect is very troubling'' (p.~99), with a note
citing \citet{DiaconisZabell1982} for how ubiquitous the effect is and
\citet{Doring1999} for its implausibility (note~15). The first prong is a claim
about what is psychologically likely and stays where introspection leaves it.
The second is a claim about effects in the world, and claims about effects are
claims about what an observer would find.

\paragraph{4. The paper answers the second prong, which is the answerable one.}
``Very troubling'' presupposes that the trouble shows up somewhere. The paper
computes where. If a population updates amnestically on level inputs, the
sequence effect has an exact observational signature, and it is lopsided: it
reaches some statistics at first order in the prior covariance $c$ and is
confined to second order in others. The debate over whether amnestic updating is
realistic thereby acquires empirical content that neither side had supplied. The
logician need not become a psychologist; he need only say which column of a
table an auditor should read.

\section*{What population statistics imply}

The classification is complete and it is settled by the two reading sequences
alone. A smooth statistic of the belief registers the sequence effect at first
order if and only if its differential at independence moves along the plane the
two sequences span; the only statistics that do not are the local dependence
measures, the believed cross-attribute association and its monotone transforms,
together with the surplus-weighted loss of the threshold decision (Propositions
PRO, ORD; Theorem LOS of the manuscript). Marginal probabilities and the
decision rate register the effect at $\Theta(c)$. The believed association
conceals it at $\bigO(c^2)$. The contrast between the two reading sequences,
$F(\PJ_{AB})-F(\PJ_{BA})$, involves neither the benchmark $\PB$ nor the mix of
sequences in the population, so any observer who records which cue an evaluator
met first can compute it. Four consequences follow for the standoff.

\paragraph{Amnesticity cannot be dismissed from dependence data.}
The statistic most naturally read as the signature of a stereotype is blind to
the sequence effect at first order. A population of amnestic updaters presents a
believed association within $\bigO(c^2)$ of the sequence-free value. Well-behaved
dependence statistics are therefore no evidence against the amnestic model, and
no amount of them settles Hawthorne's question.

\paragraph{The trouble is real but bounded.}
The same classification disciplines the attack. Amnestic updating never
manufactures cross-attribute association; the stereotype axis is untouched, and
the population's average surplus loss is second order. What is first order is
the drift of the marginals and the share of decisions that the reading sequence
decides. So the ``very troubling'' order effect has a precise shape: it lives in
levels and in the incidence of sequence-decided outcomes, not in believed
dependence and not in average welfare.

\paragraph{The model is falsifiable, cheaply.}
Amnestic updating predicts that evaluators who met the cues in opposite orders
differ at first order in levels and decisions while agreeing to second order in
association. Recording reading order and reading a level or decision statistic
is enough to test this. Conversely, while order is unrecorded or only dependence
statistics are read, no data discriminates, and Doring's demonstration cannot be
carried from his tables to any population claim. His own worked example displays
the order effect in the cells of a $2\times2$ joint over two attributes; cells
are exactly what an observer of a population does not get to read, and the step
from his tables to observable statistics is the step this paper supplies.

\paragraph{The realism complaint itself becomes a parameter.}
Hawthorne's charge was not that order effects are impossible but that we are
unlikely to ``dismiss previous experiences so completely.'' That is a claim
about a quantity: how much of the earlier impression survives. The rival
rational account of order effects, partial adjustment in the manner of
Hogarth and Einhorn, protects the \emph{first} impression; amnestic updating
protects the \emph{last}. Embedding the adjustment rule coherently makes the
two ends observable. The marginal that stays fixed as the prior covariance
varies sits on the last-read attribute under amnestic updating, on the
first-read attribute under a first-impression rule, and on neither for any
interior adjustment weight, because the last-read marginal's deviation from
its target is exactly $(1-\delta)$ times its undamped value. How completely
previous experience is dismissed is therefore not a matter for introspection.
It is a parameter, and the data locate it.

\paragraph{The blindness is not a matter of sample size.}
One clarification forestalls the obvious reply. The protected statistics are not
merely noisy: their first-order signal is zero in the population itself. The
failure is one of identification rather than estimation, and no quantity of data
overcomes it. The corresponding literature in economics, where the
informativeness of an audit is likewise settled by the structure of the
instrument rather than the size of the sample, is taken up separately in the
companion note on positioning.

\paragraph{The standoff dissolves into an instrument choice.}
Hawthorne's question, which model is the more plausible model of human agents,
was left to a psychology that was never brought to bear. The paper does not
answer it and does not need to. It replaces the question's premise: plausibility
arguments were the only currency because nobody had said what the models imply
for observation. Now the implication is a table. Which statistics one trusts,
and whether reading order is recorded, decides whether the question can be
answered at all; introspection about amnesia decides nothing.

\bigskip
\bibliographystyle{plainnat}
\begin{thebibliography}{9}
\bibitem[Diaconis \& Zabell(1982)]{DiaconisZabell1982}Diaconis, P. \& Zabell,
  S.~L. (1982). Updating subjective probability. \emph{Journal of the American
  Statistical Association} 77(380), 822--830. Read in full; formalized in
  \texttt{literature/diaconis\_zabell1982/}.
\bibitem[D\"oring(1999)]{Doring1999}D\"oring, F. (1999). Why Bayesian psychology
  is incomplete. \emph{Philosophy of Science} 66 (Proceedings), S379--S389.
  Read in full from the scanned original (Drive: \texttt{doring.pdf}); quotations
  from S379 and S388.
\bibitem[Field(1978)]{Field1978}Field, H. (1978). A note on Jeffrey
  conditionalization. \emph{Philosophy of Science} 45(3), 361--367. Read in
  full; formalized in \texttt{literature/field1978/}.
\bibitem[Garber(1980)]{Garber1980}Garber, D. (1980). Field and Jeffrey
  conditionalization (discussion). \emph{Philosophy of Science} 47(1), 142--145.
  Read in full; both numeric counterexamples reproduced exactly in
  \texttt{literature/garber1980/}.
\bibitem[Hawthorne(2004)]{Hawthorne2004}Hawthorne, J. (2004). Three models of
  sequential belief updating on uncertain evidence. \emph{Journal of
  Philosophical Logic} 33(1), 89--123. Read in full; quotations from pp.~97, 99,
  115 and note 15; both worked examples and the Update Reordering Theorem
  verified in \texttt{literature/hawthorne2004/}.
\bibitem[Jeffrey(1988)]{Jeffrey1988}Jeffrey, R.~C. (1988). Conditioning,
  kinematics, and exchangeability. In B.~Skyrms \& W.~L. Harper (eds.),
  \emph{Causation, Chance, and Credence}, vol.~1, 221--255. Dordrecht: Kluwer.
  Not read directly; its commutativity result is stated and derived as a
  corollary in Wagner (2002, Remark~3.3), and the work appears in D\"oring's
  own reference list. Cited here only for that result.
\bibitem[Wagner(2002)]{Wagner2002}Wagner, C.~G. (2002). Probability kinematics
  and commutativity. \emph{Philosophy of Science} 69(2), 266--278. Read in full
  (Drive copy is the preprint, hence quotations are located by section);
  Theorems 3.1 and 4.1 verified in \texttt{literature/wagner2002/}.
\end{thebibliography}

\end{document}
